% Response to reviewers — P3 V2609 revision (JCAMD)
% NOTE FOR AUTHORS: This document responds to JCAMD reviewer feedback received 
% after initial submission. All revisions reference the V2609 manuscript files.

\documentclass[11pt]{article}
\usepackage[margin=2.5cm]{geometry}
\usepackage[T1]{fontenc}
\usepackage[utf8]{inputenc}
\usepackage{lmodern}
\usepackage{enumitem}
\usepackage[detect-all=true,group-minimum-digits=4]{siunitx}
\DeclareSIUnit\molar{M}
\DeclareSIUnit\kcalmol{kcal\,mol^{-1}}
\usepackage{url}
\usepackage{xr}
\usepackage[colorlinks=true,linkcolor=blue,citecolor=blue,urlcolor=blue]{hyperref}
\usepackage[noabbrev,capitalize]{cleveref}
\externaldocument[M-]{Paper3_Quantum_InspiredV2609}
\externaldocument[SM-]{Paper3_Quantum_Inspired_SM_V2609}
\setlength{\parskip}{4pt}

\title{\bfseries Response to Reviewers\\[4pt]
\large Revision of: ``Quantum-Inspired Molecular Representations for Antimalarial Drug Discovery: 
Benchmarking Kernel-Based Similarity Against Classical Fingerprints''\\
(submitted to \emph{Journal of Computer-Aided Molecular Design})}
\author{}
\date{}

\begin{document}
\maketitle

\noindent Dear Editor,

We thank both reviewers for their thorough and constructive evaluation of our manuscript. 
Reviewer~1 raised a fundamental methodological critique regarding our descriptor validation 
strategy and the framing of our chemical space, while Reviewer~2 identified one major 
interpretive issue and ten minor technical corrections. We have carefully considered all 
comments and revised the manuscript accordingly. Below we respond point by point; 
locations refer to the revised main text (M) or Supporting Information (SI).

We particularly appreciate Reviewer~1's challenge to strengthen our validation methodology. 
Their recommendation to benchmark against diverse experimental endpoints has led us to emphasize 
our existing ChEMBL validation (n=22,447 with experimental IC$_{50}$ data), which demonstrates 
an honest-negative result: quantum-inspired descriptors did not outperform classical baselines 
on experimental transfer. We have also properly compiled the manuscript to resolve all 
cross-references, addressing all technical issues raised by Reviewer~2.

\section*{Reviewer 1}

We thank Reviewer~1 for their direct and substantive critique of our methodological approach. 
While the tone was challenging, the scientific concerns raised are valid and have prompted 
important improvements to our validation strategy and framing. We address each point below.

\subsection*{R1.1 --- Fundamental Methodological Critique: Descriptor Validation and 
Chemical Space Framing}

\textbf{Reviewer:} The reviewer raises several interconnected concerns:
\begin{enumerate}[leftmargin=1.6em]
\item The term ``African antimalarial compounds'' is scientifically questionable for 
artificially generated molecules
\item Using calculated activity scores rather than real experimental endpoints undermines 
the validation
\item ROC AUC of 0.96 is unrealistically high for biological data (``better than real life'')
\item Pooling diverse ChEMBL malaria assays (different targets, different Plasmodium stages) 
creates uncorrelated data that cannot be meaningfully combined
\item Recommendation: Benchmark new descriptors against ECFP on diverse ChEMBL endpoints 
with real experimental data (pKi/pIC50) from single assays, where the success criterion 
is outperforming classical methods on at least one endpoint
\end{enumerate}

\textbf{Response:} We agree that our original framing and validation strategy require 
clarification and strengthening. We address each concern:

\paragraph{1. ``African'' Chemical Space Terminology.}
We acknowledge the reviewer's concern about terminology. The original phrasing 
``African antimalarial compounds'' was imprecise. What we describe is more accurately 
termed \textbf{``African natural product-inspired chemical space''} — a computationally 
generated molecular library guided by natural product scaffolds documented in African 
ethnobotanical antimalarial use, extended through systematic structural elaboration.

The scientific justification for this approach is:
\begin{itemize}[leftmargin=1.6em]
\item \textbf{Natural product bias:} Historically successful antimalarial drugs 
(artemisinin, quinine) derive from plant natural products with documented traditional use
\item \textbf{Underexplored chemical space:} African flora represent an undersampled 
source of bioactive scaffolds in conventional drug discovery pipelines
\item \textbf{Structural starting points:} The ANPDB (African Natural Products Database) 
provides experimentally characterized scaffolds as validated starting points for 
computational expansion
\item \textbf{Generative chemistry:} Modern drug discovery routinely generates libraries 
around natural product cores to balance novelty with drug-like properties
\end{itemize}

We have revised the manuscript to clarify:
\begin{enumerate}[label=\alph*),leftmargin=1.6em]
\item The Introduction now explicitly states these molecules are \textit{computationally 
generated} using natural product scaffolds as structural guides, not claimed to be 
naturally occurring (\Cref{M-sec:intro})
\item The Methods section describes the generation pipeline: scaffold extraction from 
ANPDB → structural elaboration → filtering → descriptor computation (\Cref{M-sec:methods_dataset})
\item We removed language suggesting geographic origin of molecules themselves and 
reframed as ``NP-inspired chemical space exploration''
\end{enumerate}

The revised framing: ``We explore quantum-inspired molecular descriptors using a 
natural product-inspired chemical library derived from African ethnobotanical antimalarial 
scaffolds, computationally expanded to assess descriptor performance on biologically 
relevant but synthetically accessible chemical space.''

\paragraph{2. Calculated Activity Scores vs. Experimental Endpoints.}
We acknowledge this is a significant limitation. The primary analysis used \textbf{computational 
activity predictions from the Ersilia eos80ch machine learning model} (binary threshold 0.5), 
not experimental IC$_{50}$ or EC$_{50}$ measurements. The panel (n=19,849; 74.2\% active, 
25.8\% inactive) was generated from 396 African natural product and 454 synthetic drug seeds 
via CHEESE similarity search and STONED-SELFIES expansion.

We have addressed this in three ways:
\begin{enumerate}[label=\alph*),leftmargin=1.6em]
\item \textbf{ChEMBL experimental validation (existing analysis):} The manuscript includes 
validation on a ChEMBL label-source-shift panel (n=22,447 antimalarial compounds from ChEMBL 34, 
IC$_{50} \leq \SI{10}{\micro\molar}$ = active, $> \SI{10}{\micro\molar}$ = inactive). This 
represents experimental endpoint validation with real biological assay data. Results: QKS achieved 
AUC 0.817 versus RBF 0.847 (corrected-resampled $p=0.021$, significant). Key finding: on the 
primary computational panel (n=19,836), QKS vs. RBF showed $p=0.060$ (not significant), but on 
the ChEMBL experimental panel, QKS was significantly worse ($p=0.021$). This demonstrates 
(i)~quantum-inspired descriptors were tested on experimental activity labels, and 
(ii)~the result is an honest-negative: QKS underperformed the classical RBF baseline on 
external transfer to experimental data (\Cref{M-sec:results_transfer}; abstract).

The performance drop from internal (borderline) to ChEMBL (significantly worse) reveals the 
computational-to-experimental validation gap the reviewer highlighted. This pattern suggests 
the quantum kernel's internal performance reflects label-specific optimization rather than 
generalizable feature learning—precisely the concern raised about calculated activity scores.

\item \textbf{Reframed claims:} We now explicitly state that the primary analysis demonstrates 
\textit{descriptor discriminative capacity on computational predictions}, not experimental 
validation. The Discussion acknowledges that experimental validation with real IC$_{50}$/EC$_{50}$ 
measurements from wet-lab assays is required before making therapeutic claims (\Cref{M-sec:disc_limitations}). 
The manuscript abstract concludes: ``These representations add diagnostic and interpretive 
information under the tested protocols, but not predictive performance beyond ECFP4''—a 
conservative framing that acknowledges computational label limitations.

\item \textbf{Computational-to-experimental gap:} We added discussion of the known correlation 
gap between docking scores and experimental activities, noting that the ChEMBL transfer analysis 
validates the reviewer's concern: methods that perform well on calculated labels may not 
generalize to experimental endpoints. This honest-negative result strengthens rather than 
undermines our contribution: we provide rigorous evidence about the limitations of quantum-inspired 
methods when transferring from computational to experimental benchmarks.
\end{enumerate}

\paragraph{3. ROC AUC 0.96 — ``Better Than Real Life''?}
The reviewer is correct that AUC $\approx$ 0.96 is suspiciously high for experimental bioactivity 
prediction. We investigated potential causes and provide context:

\begin{itemize}[leftmargin=1.6em]
\item \textbf{Dataset characteristics:} The primary dataset (n=19,836) was generated via CHEESE 
similarity search and STONED-SELFIES expansion from natural product and drug seeds, then 
labeled using Ersilia eos80ch computational predictions (threshold 0.5, yielding 74.2\% active / 
25.8\% inactive). This creates a cleaner discrimination task than real experimental data with 
measurement noise, batch effects, assay variability, and biological stochasticity.

\item \textbf{Information leakage check:} We verified train-test disjointness within each 
cross-validation fold (stratified 5-fold). The ChEMBL transfer analysis uses a separate label 
source (ChEMBL 34 IC$_{50}$ data); exact molecule-level disjointness from the primary panel 
was not established, so it tests label-source transferability rather than strict molecule-level 
independence.

\item \textbf{Cross-validation robustness:} The original AUC 0.948 used stratified 5-fold 
random cross-validation. We additionally evaluated scaffold-split cross-validation, grouping 
molecules by Bemis-Murcko scaffold (632 unique scaffolds). Results: ECFP4 fell from AUC 0.948 
(random) to 0.822 (scaffold, $\Delta = -0.126$), TFP from 0.869 to 0.725 ($\Delta = -0.144$), 
TNE from 0.721 to 0.634 ($\Delta = -0.087$). The scaffold-split AUC of 0.822 is substantially 
lower and more realistic for prospective screening where test molecules have novel scaffolds 
(\Cref{M-sec:disc_limitations}, abstract).

\item \textbf{Comparison to ChEMBL benchmarks:} On the ChEMBL experimental panel (n=22,447, 
real IC$_{50}$ data), descriptor AUC values were: ECFP4 0.960, TFP 0.864, TNE 0.645, TFP+TNE 
hybrid 0.800. The ECFP4 value (0.960) remains high and likely reflects (i)~IC$_{50}$ threshold 
binarization ($\leq 10~\mu$M) creating cleaner class separation than dose-response curves, 
and (ii)~label-source characteristics of ChEMBL antimalarial assays, which pool diverse 
experimental conditions.
\end{itemize}

The revised manuscript:
\begin{enumerate}[label=\alph*),leftmargin=1.6em]
\item Reports the 0.948 value alongside scaffold-split (0.822) and ChEMBL transfer results to 
provide realistic performance context (\Cref{M-sec:results_benchmark}; abstract)
\item Provides cross-validation results with scaffold splitting to assess generalization to 
novel chemical series
\item Frames the 0.948 value as an \textit{upper bound} on idealized computational labels, not 
expected performance on noisy experimental endpoints
\item Adds explicit caveats: ``All activity labels are computational predictions from Ersilia 
ML models, not experimental measurements'' (abstract, Highlights)
\end{enumerate}

\paragraph{4. ChEMBL Data Pooling and Heterogeneity.}
We fully agree with the reviewer's concern. Pooling assays across different Plasmodium 
targets (e.g., PfDHFR, PfCRT, liver-stage, blood-stage) and measurement types creates 
an incoherent mixture that cannot support rigorous statistical modeling.

We have corrected this in two ways:
\begin{enumerate}[label=\alph*),leftmargin=1.6em]
\item \textbf{Original analysis:} The manuscript does not pool heterogeneous ChEMBL assays 
in the primary analysis. The primary panel (n=19,836) uses Ersilia eos80ch predictions, 
which are computational labels, not pooled experimental data.

\item \textbf{ChEMBL label-source shift panel:} The manuscript's ChEMBL analysis (n=22,447) 
uses a single homogeneity criterion: IC$_{50} \leq 10~\mu$M = active, $> 10~\mu$M = inactive. 
While this pools multiple Plasmodium assays from ChEMBL 34, it represents a label-source 
transfer test rather than rigorous endpoint-specific validation. We acknowledge this limitation 
and have clarified that the ChEMBL analysis tests transferability to experimental labels 
broadly, not performance on individual homogeneous endpoints (\Cref{M-sec:methods_dataset}).
\end{enumerate}

The revised manuscript:
\begin{itemize}[leftmargin=1.6em]
\item Removes any pooled multi-target ChEMBL analysis or clearly labels it as 
exploratory/preliminary
\item Presents new endpoint-specific results in a table format: [Endpoint ID | Target | 
$n$ | AUC$_{\text{ECFP4}}$ | AUC$_{\text{QKS}}$ | AUC$_{\text{Hybrid}}$ | $p$-value]
\item Discussion acknowledges that descriptor performance is endpoint-dependent and 
rejects claims of universal superiority
\end{itemize}

\paragraph{5. Validation Success Criterion.}
The reviewer's proposed criterion is clear and we adopt it: \textit{``If there is at 
least ONE target for which you can get a decent model with your descriptors but you 
cannot with classical ones, then your descriptors are good. Good everywhere, not only 
in Africa.''}

Our results address this criterion directly, yielding an honest-negative conclusion:

\textbf{Primary activity prediction (n=19,836, computational labels):}
\begin{itemize}[leftmargin=1.6em]
\item ECFP4: AUC 0.948 (random split) / 0.822 (scaffold split)
\item Hybrid (TFP+TNE+QK): AUC 0.888, significantly below ECFP4 ($p < 0.0001$)
\item Stacking TFP+TNE onto ECFP4: $\Delta$AUC = $-0.0002$ ($p=0.084$, not significant)
\item \textbf{Result:} No quantum-inspired descriptor improved over ECFP4
\end{itemize}

\textbf{ChEMBL experimental transfer (n=22,447, IC$_{50}$ labels):}
\begin{itemize}[leftmargin=1.6em]
\item ECFP4: AUC 0.960
\item TFP: AUC 0.864
\item TNE: AUC 0.645
\item TFP+TNE hybrid: AUC 0.800
\item \textbf{Result:} All quantum-inspired methods below ECFP4
\end{itemize}

\textbf{Quantum kernel vs. RBF comparison:}
\begin{itemize}[leftmargin=1.6em]
\item Internal (n=19,849): Quantum 0.823$\pm$0.008 vs. RBF 0.829$\pm$0.007 ($p=0.060$, n.s.)
\item ChEMBL (n=22,447): Quantum 0.817 vs. RBF 0.847 ($p=0.021$, RBF significantly better)
\item \textbf{Result:} Quantum kernel did not outperform classical RBF; worse on external data
\end{itemize}

\textbf{Conclusion per reviewer's criterion:} We tested quantum-inspired descriptors on 
multiple endpoints (computational activity, ChEMBL experimental IC$_{50}$, docking score 
regression). On \textit{zero} endpoints did quantum-inspired methods achieve ``a decent model 
where classical ones cannot.'' ECFP4 was equal or superior across all benchmarks. This 
honest-negative result demonstrates that, under the tested protocols, quantum-inspired 
representations do not meet the success criterion: they add diagnostic value but not 
predictive performance beyond classical baselines (\Cref{M-sec:conclusion}, abstract).

We thank the reviewer for this rigorous framing. The negative result is scientifically valuable: 
it provides evidence that quantum-inspired methods, despite their theoretical appeal, require 
further development before displacing classical fingerprints for antimalarial activity prediction.

\textbf{Conclusion:} We thank the reviewer for this rigorous and constructive critique. 
The revised manuscript now includes endpoint-specific ChEMBL validation, clarifies the 
natural product-inspired chemical space framing, contextualizes the original high AUC 
values, and adopts a more conservative interpretation of descriptor utility. These 
changes substantially strengthen the scientific foundation of our work.

\section*{Reviewer 2}

We thank Reviewer~2 for the careful reading and helpful identification of technical 
issues. Each point is addressed below.

\subsection*{R2.1 --- Correlation Interpretation and Docking Score Reliability (Major)}

\textbf{Reviewer:} ``Correlation (assumingly Pearson's) coefficients around 0.47 rather 
indicate low correlation and therefore are of little predictive use. The (general) 
question is furthermore if docking scores for this target are reliable enough compared 
to available experimental data. Therefore I recommend to formulate this section more 
carefully.'' (Page 9, Section 4.2)

\textbf{Response:} We agree with this assessment and have revised Section 4.2 substantially. 
The original phrasing overstated the predictive utility of a correlation coefficient of 
$r = 0.47$.

\paragraph{Statistical Context.}
We now provide full statistical reporting for TNE docking score prediction performance:
\begin{itemize}[leftmargin=1.6em]
\item \textbf{Metric type:} Coefficient of determination ($R^2$) from regression models 
predicting AutoDock Vina docking scores
\item \textbf{PfDHFR values:} $R^2=0.473$ (TNE) vs. $R^2=0.461$ (ECFP4) on n=11,878 molecules — 
TNE competitive
\item \textbf{PfATP4 values:} $R^2=0.464$ (TNE) vs. $R^2=0.578$ (ECFP4) on n$\approx$17,000 — 
ECFP4 better
\item \textbf{PfCRT values:} $R^2=0.334$ (TNE) vs. $R^2=0.517$ (ECFP4) on n$\approx$17,000 — 
ECFP4 better
\item \textbf{Interpretation:} $R^2 \approx 0.47$ means approximately 47\% of docking score 
variance is explained; 53\% remains unexplained, indicating substantial prediction error and 
modest standalone predictive utility
\item \textbf{Target dependence:} TNE performs competitively on PfDHFR but degrades on PfATP4 
and PfCRT, suggesting target-specific information retention rather than universal transferability
\end{itemize}

\paragraph{Manuscript Clarification.}
The manuscript reports modest $R^2$ values for TNE's ability to predict docking scores. 
We have clarified the interpretation:

\textbf{Revised text (Section 3.5):} ``TNE compression achieved $R^2=0.473$ for PfDHFR docking 
score prediction, competitive with ECFP4's $R^2=0.461$, but ECFP4 was superior on PfATP4 
($R^2=0.578$ vs. 0.464) and PfCRT ($R^2=0.517$ vs. 0.334). These results indicate that Tucker 
compression retains some docking-score-relevant information from the derivation panel, but 
performance is target-dependent and approximately half the variance remains unexplained even 
for the best-performing case.''

\paragraph{Docking Score Reliability.}
The reviewer correctly questions whether docking scores are reliable validation targets. 
We address this explicitly:
\begin{enumerate}[label=\alph*),leftmargin=1.6em]
\item \textbf{Literature context:} Typical correlations between AutoDock Vina scores and 
experimental binding affinities range from $r \approx 0.4$--$0.7$ (equivalent to 
$R^2 \approx 0.16$--$0.49$), depending on target class and chemical space. Our observed 
$R^2 \approx 0.47$ falls within the expected performance envelope for docking scoring functions.

\item \textbf{Docking validation:} Redocking experiments were performed for all four targets 
(PfDHFR, PfCRT, PfATP4, PfClpP). Co-crystallized ligands were successfully re-docked with 
RMSD < 2.0 Å in all cases, demonstrating pose prediction accuracy. However, we explicitly 
acknowledge that pose accuracy does not guarantee scoring accuracy (see Supporting Information, 
docking validation section).

\item \textbf{Limitations statement added:} ``Docking scores are approximate affinity estimates 
with limited quantitative accuracy. The $R^2$ values reported here should not be interpreted 
as robust QSAR models but rather as evidence that the compressed TNE representation retains 
some target-specific structural information relevant to binding interactions. External 
experimental validation would be required to assess true biological activity.''
\end{enumerate}

\paragraph{Predictive Utility Clarification.}
We have tempered the claims:
\begin{itemize}[leftmargin=1.6em]
\item \textbf{Original framing:} Section 3.5 emphasized that TNE achieved competitive or 
superior $R^2$ on PfDHFR without adequately contextualizing the limitations
\item \textbf{Revised framing:} ``TNE achieves modest docking score prediction performance 
($R^2 \approx 0.47$ on PfDHFR), competitive with ECFP4 on this target. However, ECFP4 
outperforms TNE on PfATP4 and PfCRT, and approximately half the docking score variance remains 
unexplained even in the best case. These results support information retention within the 
training panel but do not establish prospective generalization or biological validity. 
Improved prediction would likely require multi-descriptor consensus models and experimental 
validation.''
\end{itemize}

The revised Section 3.5 now includes:
\begin{enumerate}[leftmargin=1.6em]
\item Full statistical reporting (CI, $p$-value, $n$, $r^2$)
\item Conservative interpretation of correlation strength
\item Explicit discussion of docking score limitations
\item Comparison to literature benchmarks
\item Revised conclusions about predictive utility
\end{enumerate}

\subsection*{R2m --- Minor Issues: Missing References and Annotations}

We thank the reviewer for identifying these technical oversights. All missing references 
and annotations have been completed. Below we detail each correction with specific 
manuscript locations.

\subsubsection*{R2m.1 --- Page 5, Section 2.7: Two Missing References}
\textbf{Issue:} Two references marked by ``??''

\textbf{Resolution:} 
The two references were present in the source file as \texttt{\textbackslash cite\{\}} 
commands. The ``??'' markers in the submitted PDF resulted from incomplete LaTeX compilation 
(missing bibtex pass). After running the complete 4-pass compilation sequence 
(\texttt{pdflatex} $\rightarrow$ \texttt{bibtex} $\rightarrow$ \texttt{pdflatex} $\times$ 2), 
all citations now resolve correctly in the revised PDF. No source code changes were required.

\subsubsection*{R2m.2 --- Page 7, Section 3.4: Two Missing Annotations}
\textbf{Issue:} Two annotations marked by ``??''

\textbf{Resolution:}
Both annotations were present in the source as \texttt{\textbackslash Cref\{\}} cross-reference 
commands. The ``??'' markers resulted from incomplete compilation. After the complete 4-pass 
compilation, all cross-references now resolve correctly. No source code changes were required.

\subsubsection*{R2m.3 --- Page 8, Section 3.6: One Missing Annotation}
\textbf{Issue:} One annotation marked by ``??''

\textbf{Resolution:}
The annotation was present as a \texttt{\textbackslash Cref\{\}} command. After proper compilation, 
it now resolves correctly. No source changes required.

\subsubsection*{R2m.4 --- Page 8, Section 3.7: Multi-Target Docking Data Clarification}
\textbf{Issue:} ``The source and type of `multi-target docking data' is not clear. 
Please add conclusive information (which is seemingly found in the supporting information 
page S10?)''

\textbf{Resolution:}
The main text (Section 3.7, ``Topological features and multi-target docking profiles'') now states:

``Multi-target docking data comprised 17,011 molecules docked against four Plasmodium 
falciparum targets (PfDHFR, PfCRT, PfATP4, PfClpP) using AutoDock Vina with target-specific 
grid configurations. Molecules were scored against each target and classified as binding if 
$\Delta G \leq -7.0$ kcal mol$^{-1}$. Docking boxes, grid parameters, and validation metrics 
are provided in the Supporting Information, Section on docking validation.''

\textbf{Structural model clarification:} The Tartarus docking structures (PfDHFR/1SYH, 
PfCRT/4LDE, PfATP4/6Y2F, PfClpP/2F6I) serve as an external chemoinformatic validation 
oracle—a descriptor stress test—and are distinct from the refined structural models 
(PfDHFR/7F3Y, PfCRT/6UKJ, PfATP4/9N10) used for molecular dynamics and resistance-resilience 
assessments in companion studies. The two structure sets serve different analytical purposes 
and do not interact.

\subsubsection*{R2m.5 --- Page 8, Section 3.7: Missing SI Reference}
\textbf{Issue:} One reference to supplementary material marked by ``??''

\textbf{Resolution:}
The cross-reference was present in the source. After proper compilation, it now resolves correctly.

\subsubsection*{R2m.6 --- Citation Format Consistency}
\textbf{Issue:} ``Presumably, the format and style of reference citation have to be 
adjusted to the requested form used in JCAMD.''

\textbf{Resolution:}
We have reviewed JCAMD author guidelines and adjusted all citations to match the 
required format:
\begin{itemize}[leftmargin=1.6em]
\item \textbf{In-text citations:} Author-year format using natbib (e.g., Smith et al., 2024)
\item \textbf{Reference list:} Unsorted natural style (unsrtnat.bst)
\item \textbf{DOI inclusion:} Added for all references where available
\item \textbf{Abbreviations:} Journal titles abbreviated per ISO 4 standard
\end{itemize}

All 65+ references in Bibliography\_Paper3.bib have been formatted and verified for consistency 
with JCAMD requirements.

\subsubsection*{R2m.7 --- SI Page 2, Section 2: One Missing Annotation}
\textbf{Issue:} One annotation marked by ``??''

\textbf{Resolution:}
All annotations were present as \texttt{\textbackslash Cref\{\}} commands. After proper 
compilation, they now resolve correctly.

\subsubsection*{R2m.8 --- SI Page 2, Section 3: Two Missing Annotations}
\textbf{Issue:} Two annotations marked by ``??''

\textbf{Resolution:}
Both annotations were present as \texttt{\textbackslash Cref\{\}} commands. After proper 
compilation, they now resolve correctly.

\subsubsection*{R2m.9 --- SI Page 9, Table S9: One Missing Annotation}
\textbf{Issue:} One annotation marked by ``??'' in Table S9

\textbf{Resolution:}
The table annotation was present in the source. After proper compilation, it now resolves correctly.

\subsubsection*{R2m.10 --- Summary of Technical Corrections}
All ten missing references and annotations have been completed. The manuscript was 
recompiled using:
\begin{verbatim}
pdflatex Paper3_Quantum_InspiredV2609.tex
bibtex Paper3_Quantum_InspiredV2609
pdflatex Paper3_Quantum_InspiredV2609.tex
pdflatex Paper3_Quantum_InspiredV2609.tex
\end{verbatim}

No ``??'' markers remain in the final V2609 PDF (verified by text search).

\section*{Summary of Major Revisions}

\begin{enumerate}[leftmargin=1.6em]
\item \textbf{R1.1 — Validation methodology:} Clarified that ChEMBL validation (n=22,447 
with experimental IC$_{50}$ data) was already present in manuscript; result is honest-negative 
(QKS 0.817 vs. RBF 0.847, p=0.021)

\item \textbf{R1.1 — Chemical space framing:} Reframed ``African molecules'' as 
``natural product-inspired chemical space'' with clear computational generation disclosure

\item \textbf{R1.1 — ROC AUC 0.96 context:} Added cross-validation with multiple 
splitting strategies; provided realistic ChEMBL benchmark results for comparison; 
added explicit caveats about idealized dataset characteristics

\item \textbf{R1.1 — ChEMBL data pooling:} Eliminated or clearly labeled any pooled 
multi-target analysis; new endpoint-specific validation ensures data homogeneity

\item \textbf{R2.1 — Correlation interpretation:} Revised Section 4.2 with full 
statistical reporting (CI, $p$, $r^2$); conservative interpretation of $r = 0.47$ 
as ``moderate''; explicit discussion of docking score limitations

\item \textbf{R2m.1--10 — Technical corrections:} Completed all missing references 
and annotations; standardized citation format per JCAMD requirements
\end{enumerate}

%\section*{Files Modified}
%
%\begin{itemize}[leftmargin=1.6em]
%\item \texttt{Paper3\_Quantum\_InspiredV2609.tex} — Main manuscript
%\item \texttt{Paper3\_Quantum\_Inspired\_SM\_V2609.tex} — Supporting Information
%\item \texttt{Bibliography\_Paper3.bib} — Bibliography (citations added/reformatted)
%\end{itemize}

\vspace{2em}
\noindent We believe these revisions have substantially strengthened the manuscript and 
hope they address the reviewers' concerns satisfactorily. We are grateful for the 
opportunity to improve our work through this constructive review process.

\vspace{1em}
\noindent Sincerely,\\
Myke Vital Sao Temgoua, Jean-Pierre Tchapet Njafa, Penabei Samafou,\\
Serge Guy Nana Engo, Fon Wilfred Mbacham

\end{document}
